\documentclass[11pt,a4paper]{article}
\usepackage{iftex}
\ifPDFTeX
  \usepackage[T1]{fontenc}
  \usepackage[utf8]{inputenc}
  \usepackage{lmodern}
  \input{glyphtounicode}
  \pdfgentounicode=1
\else
  \usepackage{fontspec}
  % Kerning off: kerned pairs ("A WS", "T ree") split words for ATS text extraction.
  \setmainfont{lmroman10-regular}[Extension=.otf, BoldFont=lmroman10-bold,
    ItalicFont=lmroman10-italic, BoldItalicFont=lmroman10-bolditalic, Kerning=Off]
\fi
\usepackage[margin=0.9in]{geometry}
\usepackage[hidelinks]{hyperref}
\hypersetup{pdftitle={Abhinav Kaushik - Cover Letter - openai}, pdfauthor={Abhinav Kaushik}}

\pagestyle{empty}
\setlength{\parindent}{0pt}
\setlength{\parskip}{10pt}
\raggedright

\begin{document}

{\Large\bfseries Abhinav Kaushik}\\
Applied AI Engineer\\
New Delhi, India \textbar{} +91 8700571859 \textbar{} \href{mailto:aabhinav.kaushik@gmail.com}{aabhinav.kaushik@gmail.com} \textbar{} \href{https://www.linkedin.com/in/abhinav-kaushik10/}{LinkedIn}  

September 27, 2026

Hiring Team\\
openai

\textbf{Re: Applied AI Engineer, Startups (Codex)}

Dear Hiring Team,

I am applying for the Applied AI Engineer role on the Startups team, focused on Codex. Most of my work has been building practical AI applications quickly with the people who will use them, then making them reliable enough to stay in use. Turning a promising idea into a workflow a team keeps using is what ties my work together.

Codex is changing how small teams write and ship software, and helping founders get real, lasting value from it is work I would enjoy both in code and in front of a room. My fiancée lives in Berlin, and we plan to move to London together and settle there long term.

I would bring a builder's view to the role. At Viavi I prototyped 15+ AI applications from client requirements, and 7 of them moved into production or paid client work. At United Health Group I shipped a lab-extraction agent that automated the work behind 800 Jira tickets, with evaluation pipelines that measured its reliability.

Outside work I am building ThinkTree.AI, a non-profit AI learning platform used by NGO teachers to support students with ADHD. It keeps me focused on whether a tool actually helps the people using it, which is the same question I would ask of every startup workflow. I would welcome a conversation.

Sincerely,\\[4pt]
Abhinav Kaushik

\end{document}